\label{sec:discussion}

\subsection{Critical evaluation of the design} % REQUIRED HEADING

This subsection interprets the results of Section~\ref{sec:results}, evaluates the design choices in hindsight and identifies what must change before the requirements can be considered proven.
\label{sec:dis-critical}

Three limitations determine every conclusion below.
The first is one of platform: the reproduction and classification timings of qualification tests~3 and~4 were taken on the ODROID~N2+, which is the target platform, while the accuracy measurements of qualification tests~1 and~2 were taken on the development computer, where the platform affects only how long a measurement took and not what it was.
The second is that the evidence on synthesis quality is partly circular, because the classifier that selects the best of $N$ candidates belongs to the same family as the one that reports the score.
The third is economic, and it is set out under economic feasibility in Section~\ref{sec:dis-crit}: the processing platform was constrained by a \mbox{R5\,000} budget, and the consequences of that constraint fall almost entirely on the one requirement that was not met.
Where a figure is affected by any of the three, this is stated with the figure.

    \subsubsection{Interpretation of results}
    \label{sec:dis-interp}

        \textbf{R2, writer identification: the strongest result in the report, and the one most nearly separated from its target.}
        The writer was named correctly on 153 of 156 lines held out from the classifier, 98.08\%, and on all 11 writers when each writer's held-out lines are pooled by ink.
        The point estimate clears the 95\% target by 3.08 points, but the honest statement is weaker than that: the Wilson interval runs from 94.5 to 99.3\%, so its lower bound sits 0.5 points \emph{below} the target and 156 lines are not quite enough to separate the result from the specification at 95\% confidence.
        Roughly 250 held-out lines at the same error rate would be.
        What the pooled result adds is that no writer was misnamed as a whole: the three errors are isolated lines on three different writers, so a page-level decision, which is what the system actually makes, was right every time.

        \textbf{Why the writer classifier that reads real ink is deliberately not the best one available.}
        An ink-based classifier scored higher than the stroke-normalised one on photographed handwriting, and it was nevertheless replaced.
        The reason is that ink density and stroke thickness are among the strongest cues available in a photograph and are exactly the cues a constant-width ballpoint in a plotter cannot vary: a classifier that leans on them cannot judge the machine's own output, which is the measurement R1 depends on.
        Redrawing every crop at one constant pen width before classifying discards those cues deliberately, costs accuracy on photographs, and buys an instrument that works on both real ink and drawn output.
        This is the project's decision heuristic in its clearest form: where two options differed, the one chosen was the one that survived a change of pen, camera or lighting, even at a cost in headline accuracy.
        The normalisation width itself was not guessed but swept, and 0.17 x-heights was selected on held-out accuracy.

        \textbf{R2, text: 94.39\% against a 95\% target, and the shortfall is 0.61 points.}
        The text half was measured on 1\,076 lines held out from the reader, of which 976 are the public corpus's own validation partition and 100 are the photographed writers' seeded split.
        Two features of that measurement matter more than the headline.
        First, the two sets agree to 0.03 points, 94.39\% against 94.42\%, although one is a public benchmark of many writers cut by its own distributor and the other is five writers photographed with this project's webcam and cut by this project's segmenter.
        That agreement is the strongest available evidence that the reader generalises rather than having learned the particular conditions of the project's own pages, and it is also evidence that the segmenter is not systematically damaging the crops it produces.
        Second, the figure is a measurement on genuinely unseen material: the transcript lookup reported in qualification test~2 established that three of the database writers' pages are inside the reader's training corpus, and those pages were excluded rather than quoted.
        An earlier draft of this report quoted 91.02\% over 155 lines, a figure that included contaminated material; the present figure is both higher and sound, which is the opposite of the usual direction and is worth stating plainly.

        \textbf{Why the text half falls short, and what the gap is made of.}
        Three causes are identifiable and none of them is a defect in the recogniser itself.
        The first is the segmenter: when two written lines sit close together it occasionally carries a descender or a stray letter from the line above into the crop below, and the reader is then charged with a character that was never part of that line.
        This is infrequent, but it sets a ceiling on page-to-text accuracy that no amount of training can lift.
        The second is the deliberate absence of a language model.
        The gap between 94.39\% character accuracy and 73.47\% word accuracy localises the residual error: one wrong character spoils an entire word, so the errors are concentrated in a minority of words rather than spread thinly, and that is precisely the error structure a lexicon or an $n$-gram decoder corrects.
        A dictionary-constrained decoder was not used, and the reason is a requirement rather than an oversight: R1 asks the system to reproduce \emph{arbitrary} digital text, including proper names and strings that are not words, so constraining the decoder to a lexicon would raise the measured accuracy on benchmark prose while making the system worse at the task it is specified to do.
        The third is training budget, discussed under economic feasibility below.
        Of the three, only the first is a fault; the second is a defensible trade and the third is a constraint.

        \textbf{R1, representation fidelity: 84.41\% against 85\%, and the shortfall is concentrated in a measurable, explainable place.}
        Over the 184 lines the robot itself drew, the text was read back at 84.41\% character accuracy, 0.59 points below the target, while the stylistic half was met: the writer was identified from the machine's own output on 86.4\% of lines and on all 14 pages when a page's lines are pooled.
        The informative figure is not the mean but the spread: character accuracy by writer runs from 60.97\% to 97.01\%, a range of 36 points, which is far too wide to be noise and points at a property of the writers rather than of the machine.
        That property was measured. Ranking the eleven writers by connected ink components per character and by character accuracy gives a Spearman correlation of $+0.845$, with the four writers above 93\% being the four least joined hands and two of the three lowest being the two most joined.
        The same ranking against writer identification gives $-0.027$, that is, nothing.
        So cursiveness governs how legibly the machine reproduces a hand and has no bearing on whether the hand remains recognisable, which is a coherent result rather than a coincidence.

        \textbf{Why cursiveness costs legibility, and why that is a property of the chosen method.}
        The synthesiser composes a line from glyph paths traced out of the writer's own ink and then joins them.
        For a printed hand, the isolated letterform harvested from the page is very close to the letterform the writer actually produces in context, so the reconstruction is faithful.
        For a joined hand it is not: the writer never draws those letters in isolation, every letter is deformed by its neighbours, and a library of isolated shapes is therefore a worse model of a cursive writer than of a printed one.
        The shortfall on R1 is thus the predictable cost of a design decision taken for other reasons.
        A glyph library was chosen over a generative image model because it preserves the writer's actual letterforms, needs no accelerator to run on the board, and can be inspected and corrected glyph by glyph; a diffusion or adversarial synthesiser would model context directly and would be expected to do better on exactly the cursive cases, at the cost of all three of those properties and of a platform the budget did not allow.
        Four writers clearing 93\% shows the method works where its assumption holds.

        \textbf{The stylistic half of R1 is met on one measure and missed on the measure the proposal actually specifies.}
        The proposal asks for 85\% stylistic correlation. Measured as a correlation, over nine geometric features between each writer's real lines and the system's reproduction of the same text, the pooled coefficient is 0.714 against 0.85, so the specification is not met and the honest verdict on R1's stylistic half is a failure by 0.136.
        Measured as writer identification from the machine's own output it is met, at 86.4\% of lines and 14 of 14 pages.
        Both numbers are reported because they answer different questions, and the specification's own question is the one that fails.

        The per-feature breakdown says where the fidelity is lost, and it is coherent with the R1 character result.
        Slant correlates at 0.996, ascender extent at 0.974 and character width at 0.936, so the \emph{letterforms and their proportions} transfer almost exactly, which is expected because the synthesiser draws the writer's own traced paths.
        Word gap correlates at 0.605, the ascender-to-descender ratio at 0.509 and descender extent at only 0.148, so what fails to transfer is the \emph{spacing and the vertical placement between letters} rather than the letters themselves.
        That is the same defect the character accuracy showed from the other direction, and it localises the remaining work: the glyph library is adequate and the layout model that places those glyphs is not.

        Two caveats remain on the writer-identification figure and neither applies to the correlation.
        The classifier that reports the score also selects the best of $N$ candidates, so a best-of-$N$ search inflates the measure it is judged by; qualification test~1 limits this by classifying the G-code before any page was drawn, which removes the drawing from the loop but not the classifier.
        The overlay test gives an area under the curve of 0.71 against 0.5 for output carrying no writer information at all, which is independent of the classifier and indicates real but partial fidelity, consistent with a correlation of 0.714 rather than with the identification figure.
        A human panel was not carried out.

        \textbf{R3, reproduction time: met on the target platform, with the margin in an unexpected place.}
        All 36 complete reproductions in the log met the 4~s per character allowance, and they were measured on the ODROID~N2+ driving the real gantry, not on the development computer, so this is an on-target result.
        The margin, however, is not uniform: the worst case used 96.9\% of its allowance and the best 33.6\%, and the worst case is the \emph{shortest} block, not the longest.
        The reason is a mismatch between how the cost scales and how the allowance scales.
        Drawing is a per-character cost and tracks the allowance, but the best-of-$N$ search is paid once per block, so its contribution per character falls as the block lengthens: in the worst run at each length, synthesis cost 0.617~s per character on a 43-character block and 0.156~s per character on the 378-character paragraph.
        The allowance, being strictly proportional to length, gives a short block almost no room to pay for a fixed search.
        Taking the measured synthesis cost of the shortest block and the measured drawing rate, the allowance is exhausted at about 35 characters, which means R3 is met over the range tested and would be expected to fail below roughly that length.
        This is recorded under expected failure conditions rather than hidden in the margin.

        \textbf{R4, classification time: recognition passes, segmentation fails badly.}
        Recognition ran at 8.86~ms per character against the 25~ms allowance, a margin of about 2.8 times, so the neural component is not the problem.
        Page segmentation took 35 to 57~s against the roughly 2~s the equation allows for capture and conditioning, a shortfall of more than an order of magnitude.
        Since $x$ appears only in the recognition term, a longer page does not buy proportionally more time, so the segmenter's per-page cost cannot be amortised.
        The profile shows the cause: line grouping and the chain search dominate, at about 70\% of the time, and are written in pure Python.
        Moving two inner routines to C reduced the total by 25\% to 32\%, but the grouping still dominates: accelerating the wrong part of a program cannot fix it.

        \textbf{R5, positional accuracy: met on the specimen measured, and the residual error is explained rather than merely bounded.}
        The worst of the four sides of a 100~mm square was 0.30~mm from nominal and the mean absolute error 0.175~mm, both inside the 0.5~mm limit.
        What makes this more than a single lucky measurement is that the triangle explains where the error comes from.
        A square cannot detect a squareness error, because a sheared square still has four equal sides; the triangle can, because its two sloped sides are equal only if the axes are perpendicular.
        Those two sides differed by 0.70~mm although they are equal by construction, and solving the shear from them gives a single apex offset that predicts both measured sides to within 0.02~mm, which is the resolution of the caliper.
        One unknown fitted to two measurements with residuals at the reading resolution is a validated model, not a curve fit, and it says the dominant residual is a $0.706^{\circ}$ non-perpendicularity of the axes rather than a scale error or lost steps.
        That is a mechanical alignment fault, correctable by squaring the frame, and it is the most useful thing the test produced.
        The limitation is equally plain: one square and one triangle support a compliance statement and a diagnosis, and do not support a repeatability figure or a tolerance.

        \textbf{R6, stored styles: met, and the margin is one style.}
        Eleven writer styles are trained and permanently stored against a minimum of ten, each with a glyph library and a style profile, and the classifier checkpoint records the same eleven names.
        The set was reduced from thirteen during this work by withdrawing two writers and retraining, which is a reduction of the margin from three styles to one and is reported as such.

        \textbf{Requirements met, stated plainly.}
        R3, R5 and R6 are met, and both field conditions are met.
        R2 is met on writer identification, at 98.08\% against 95\%, and falls 0.61 points short on text.
        R1 falls 0.59 points short on its character half and, on the correlation its specification names, 0.136 short on its stylistic half; it is met only on the writer-identification reading of that half.
        R4 is not met, by a factor of about 2.4 on the per-character term and about 8 on the fixed term.
        So three of the six core requirements are met outright, two are met in part and miss by less than a point on the half that is close, and one fails by a wide margin. Every requirement was measured; none is left unassessed.

        \textbf{Where this system sits against published work.}
        Table~\ref{tab:dis-compare} places the measured figures of Section~\ref{sec:results} beside the published systems surveyed in Section~\ref{sec:literature}.
        One row of that table is a genuine comparison and the rest are not, and the difference matters.
        The reader's 94.39\% character accuracy was measured on 976 line crops of the public IAM corpus as distributed by Teklia, cut by its distributor and held out by its own partition, which is the same corpus the two published recognition rows were obtained on.
        On that corpus a recogniser written from first principles, trained on one machine without a framework supplying any layer, gradient or optimiser, reads 12.6 points above the bidirectional recurrent network of Graves \emph{et al.} and 2.0 points below the CNN--BiLSTM--CTC of K{\i}z{\i}l{\i}rmak \emph{et al.}, whose layer pattern this implementation follows.
        That is the single most useful calibration in the report: it places the weakest of the six requirements against the state of the art rather than against a target, and it shows the 0.61-point shortfall on R2 to be a shortfall against a specification and not against what the method can do.
        Two differences still favour this implementation and must be stated: the validation partition used here is the corpus's own and is not necessarily the same partition as the test sets the published figures were obtained on, and the published works report on the full IAM test material with research-scale training behind them.
        The synthesis rows are not comparable at all, being Fr\'echet inception distances, a distance between image distributions for which lower is better and for which no accuracy equivalent exists; they are included to show the class of method, and are comparable only to each other.

\begin{table}[!htbp]
\centering
\fontsize{8}{9.4}\selectfont
\setlength{\tabcolsep}{4pt}
\begin{tabular}{|p{3.0cm}|p{4.0cm}|p{3.0cm}|p{3.2cm}|}
\hline
\textbf{Work} & \textbf{Method} & \textbf{Data} & \textbf{Reported figure, with its metric} \\
\hline
Graves \emph{et al.}~\cite[p.~855]{graves2009} & Bidirectional recurrent network, CTC output & IAM lines, offline & 74.1\,\% word, 81.8\,\% character accuracy \\
\hline
K{\i}z{\i}l{\i}rmak \emph{et al.}~\cite[p.~1]{kizilirmak2023} & CNN--BiLSTM--CTC, the layer pattern this project follows & IAM lines & 3.59\,\% CER, 9.44\,\% WER (that is 96.41\,\% character accuracy) \\
\hline
Kang \emph{et al.}~\cite[Tab.~1]{kang2020} & Adversarial image synthesis & IAM & FID 120.1 (seen styles) to 130.7 (unseen, out-of-vocabulary); real data scores 90.4 \\
\hline
Dai \emph{et al.}~\cite[Tab.~3]{dai2024} & Diffusion-based synthesis, one reference image & IAM, unseen styles, out-of-vocabulary & FID 102.8 \\
\hline
\textbf{This implementation}, text recognition & CNN--BiLSTM--CTC written from first principles, no framework supplying a layer, gradient or optimiser & \textbf{IAM lines, 976 held-out crops of the distributor's own validation partition} & \textbf{94.39\% character, 73.67\% word} \\
\hline
\textbf{This implementation}, text recognition & The same checkpoint, on this project's own photographed pages & 100 held-out lines, 5 writers, webcam capture & 94.42\% character, 70.74\% word \\
\hline
\textbf{This implementation}, writer identification & The same backbone on stroke-normalised input, 11-class closed set & 156 held-out lines, 11 writers & 98.08\% per line, 11 of 11 writers pooled \\
\hline
\textbf{This implementation}, synthesis read back & Glyph-library synthesis from the writer's own traced letters & 184 lines the robot drew, 11 writers & 84.41\% character, writer identified on 86.4\% of lines \\
\hline
\end{tabular}
\caption{The measured results of this project beside the published systems surveyed in Section~\ref{sec:literature}. Only the first of this implementation's rows is a like-for-like comparison, being measured on the same public corpus as the two recognition rows above it; the writer-identification and synthesis rows have no published counterpart on the same material and are given for completeness. The two synthesis rows above are Fr\'echet inception distances and are comparable only to each other.}
\label{tab:dis-compare}
\end{table}


    \subsubsection{Critical evaluation}
    \label{sec:dis-crit}

        \textbf{Economic feasibility, and what the budget decided.}
        The project was built to a \mbox{R5\,000} budget, and the processing platform was the item that constraint bound hardest.
        The ODROID~N2+ was obtained second hand at no cost, which is what kept the build inside the budget once the gantry mechanics, motors, drivers, supply and camera were accounted for.
        Two consequences followed from that choice and neither was incidental.
        The first was practical: the board has no HDMI output, so all development, testing and demonstration had to be done headless over the network, and the browser interface described in Section~\ref{sec:system} is not merely a convenience but the only means of operating the machine.
        The second is the direct cause of the one requirement that was not met.
        Boards with substantially more processing capability were available and were considered, and were rejected on price: a Rock~5B or a Jetson Nano would have cost a large fraction of the whole budget on its own, against a board that cost nothing.
        No measurement of either alternative was made, so this report offers no figure for how much faster they would have been; what it can state is the size of the deficit that has to be covered, which Section~\ref{sec:dis-unsolved} puts at a factor of about 2.4 on the network stages and about 8 on the per-page image stages.
        The honest summary is that R4 was traded for the rest of the system: within a \mbox{R5\,000} budget the money went into a machine that draws accurately and a platform that could run the algorithms at all, and the classification-time specification was the casualty.

        \textbf{Writing the networks from first principles, and what that cost.}
        Every layer, gradient and optimiser used in training and inference was written out rather than taken from a framework, which is the project's largest single claim to original work and is also measurable as a cost.
        Loading the same weights into equivalent library modules and running each page as one batched forward pass, rather than the hand-written forward pass one crop at a time, is faster by a factor of about 2.8 on the same machine and on the same pages, with the text identical.
        That factor is not a criticism of the hand-written code, which has to be correct and legible rather than fast; it is the price of the design decision, and it happens to be most of what R4 needs.
        In hindsight the right structure would have been to keep the from-first-principles implementation as the reference and the thing that earns credit, and to batch its forward pass per page rather than per crop, which is an implementation change of the same code and not a move to a library.
        That was not done, and it is the clearest piece of available performance left on the table.

        \textbf{Training by the two routes is not equivalent, and the difference decided which checkpoint was used.}
        The writer classifier was trained twice on the same data, the same architecture and the same stroke normalisation at 0.17 x-heights: once through the library route and once through the from-first-principles trainer.
        Both were then measured on the two populations the classifier has to serve, and Table~\ref{tab:dis-routes} gives the result.
        The from-first-principles checkpoint meets R2 on held-out real handwriting, at 95.51\% against the 95\% target, so the claim that the training code works is supported by a checkpoint that satisfies the requirement.
        It is nevertheless not the checkpoint the measurements of Section~\ref{sec:results} were taken with, and the reason is the second population.
        On the pages the robot drew, where the classifier is the measuring instrument for R1 rather than the subject, it reaches 80.54\% against the library-trained checkpoint's 85.95\%, and R1's stylistic target is 85\%.
        Installing it would therefore convert R1's writer-identification half from met to not met and would lose one of the eleven pages when lines are pooled.
        The decision was to measure with the better instrument and to report the weaker one as evidence about the training route, which is what it is evidence of; the alternative, installing a checkpoint that is 5.41 points worse at the job it has to do in order to make a sentence about libraries simpler, would be choosing the claim over the measurement.

\begin{table}[h]
\centering
\fontsize{10}{12}\selectfont
\begin{tabular}{|l|r|r|r|}
\hline
\textbf{Population} & \textbf{Library route} & \textbf{First principles} & \textbf{Difference} \\
\hline
Held-out real lines (R2), per line   & 98.08\,\% & 95.51\,\% & $-2.56$ \\
\quad 95\,\% Wilson interval          & [94.5, 99.3] & [91.0, 97.8] & --- \\
\quad writers named when pooled       & 11/11 & 11/11 & --- \\
\hline
Lines the robot drew (R1), per line  & 85.95\,\% & 80.54\,\% & $-5.41$ \\
\quad pages named when pooled         & 11/11 & 10/11 & $-1$ \\
\hline
\end{tabular}
\caption{The same writer classifier trained by the two implementation routes,
measured on 156 held-out lines of real handwriting and 185 lines the robot
drew. The from-first-principles checkpoint meets the 95\,\% target of R2 and
falls below the 85\,\% stylistic target of R1, which is why the library-trained
checkpoint was used as the measuring instrument.}
\label{tab:dis-routes}
\end{table}

        \textbf{Line-level reading instead of character separation.}
        The design replaced the proposal's character separation by separation into lines and a sequence recogniser.
        This was right for connected handwriting, where character boundaries do not exist, and it makes the accuracy independent of a separation error; the price is that a whole line must be read by one network, at a cost per character higher than the 25~ms the proposal assumed.

        \textbf{Best-of-$N$ search with the judges inside the loop.}
        The search with a recogniser and a writer classifier as judges was good for quality, because it prevents the synthesis from drawing an illegible or wrong-style sample.
        In hindsight it was poor for verification, because it removed the independence of the evaluation, and for time, because each of the $N$ candidates costs a rendering and two network passes.
        The qualification protocols address the first, and a lower $N$ or a cheaper first-stage filter the second.

        \textbf{Pure Python for the segmentation, and the cheap fix that does not work.}
        The proposal required that the processing run on the embedded platform and be written without library functions for the algorithm itself.
        Writing the image processing in Python made the design transparent and correct, and it is the reason the per-page image stages cost 16.1 to 16.8~s on the board against the 2~s the specification allows for everything that does not scale with the character count.
        The obvious remedy is to work at lower resolution, since the filters and the connected-component labelling scale with the pixel count, and halving each axis should be a factor of four.
        It was measured rather than assumed, and it does not work.
        Halving the linear resolution did reduce the image stages by a factor of 3.4 to 7.9, but the segmenter found fewer lines on every page tested, the text read back from the reduced page agreed with the full-resolution reading on only 34 to 45\% of its characters, and on one page the segmenter found no line centres at all and raised an exception instead of returning an empty result.
        A speed-up that loses lines is not a speed-up, so this route is closed: the line grouping has to be redesigned rather than fed a smaller picture.
        By Amdahl's law, if a fraction $p$ of the time is accelerated by a factor $s$, the total speed-up is $1/((1-p)+p/s)$, so accelerating anything other than the grouping cannot exceed about $3.3$ even if the accelerated part became free.

        \textbf{Open-loop gantry with step counting.}
        Using the stepper motors without encoders is a standard, low-cost choice, and the budget made it the only choice; its accuracy depends on missed steps, and no sensor reveals them.
        The measurement vindicates it: the dominant residual on the drawn geometry is a $0.706^{\circ}$ non-perpendicularity of the two axes, which is a frame alignment error and not a stepping error, and the two axis scales after calibration are $+0.250\%$ and $+0.000\%$.
        In other words the open-loop drive is not what limits R5, and the margin available would not be improved by encoders until the frame is squared.
        Running the driver without acceleration ramps, however, remains a real defect: the generator models an acceleration of 400~mm/s$^2$ that the driver does not execute, so the real motion differs from the model and a lost step is more likely at the start of a stroke at higher feed.


    \subsubsection{Unsolved problems}
    \label{sec:dis-unsolved}

        The following problems were not solved.

        \begin{enumerate}
          \item \textbf{The configuration measured is not the configuration deployed.} This is the most consequential open item, because it affects how every figure in Section~\ref{sec:results} should be read.
          The accuracy measurements were taken with the reader checkpoint \texttt{paper\_cnn\_bilstm\_ctc\_scratch\_best.npz} and the writer classifier \texttt{author\_shape\_11\_pw170\_weights.npz}, the latter applied with stroke normalisation at 0.17 x-heights because that is what it was trained on.
          The server as it stands selects a different reader, because \texttt{paper\_cnn\_bilstm\_ctc\_joint\_best.npz} stands first in its candidate list, and it sets \texttt{PAGE\_NAMER\_STROKE\_NORM = False} while using the stroke-normalised classifier, so it feeds un-normalised crops to a network trained on normalised ones.
          The deployed reader was separately measured on the same held-out material and reads 94.06\% against 94.39\%, so the reader substitution costs 0.33 points; the classifier mismatch is a train-and-serve disagreement whose cost was not measured and may be larger.
          Neither is a design fault and both are one-line configuration changes, but until they are made the report's accuracy figures describe a configuration the server does not currently load.
          \item \textbf{Training data that the corrected labels have not reached.} Seven pages of one database writer are currently skipped when the data set is built, because the cached crops were written before the transcripts were corrected and the two now disagree on the line count.
          Those pages are therefore absent from training although their labels are now right, and the caches must be rebuilt before the next training run so that the corrected material is actually used.
          \item \textbf{R4 not met, and the shortfall is in both halves, not only the segmentation.}
          An earlier reading of this attributed the whole deficit to segmentation. The on-target measurement contradicts it: on the ODROID~N2+ the two network stages account for 67.6\% of the classify time and the four image stages for 32.4\%.
          Because the allowance has only one term that grows with the page, the two halves must be judged separately, and both fail.
          The networks cost 0.0394 to 0.0602~s per character against the 0.025~s per character term, so they exceed their term by a factor of 1.6 to 2.4 \emph{at every page length}; the image stages cost 16.1 to 16.8~s against the fixed 2~s term, a factor of 8.1 to 8.4.
          \item \textbf{Two latent faults in the segmenter.} Both were found by driving it outside its normal operating point and neither is reached by the deployed pipeline today.
          A page that yields no line centres raises an \texttt{IndexError} in the seam-cut construction instead of returning an empty result, which is what the half-resolution experiment struck.
          Separately, two pages of one writer each lose a written line because the band in question is assigned no ink at all by the ownership step, so the crop builder has nothing to cut; the transcripts for those pages are correct, and the fault is in ownership rather than in the ink threshold or the labels.
          \item \textbf{Optimistic validation of the photographed writers.} Their held-out material is a seeded 15\% of lines drawn from within their own pages rather than whole pages set aside, so a held-out line can share a page with training lines.
          The public-corpus measurement does not have this weakness and agrees with the photographed figure to 0.03 points, which is reassuring but is not the same as fixing the split.
          \item \textbf{Circular synthesis evaluation.} No independent judge or human panel was used, so the stylistic half of R1 rests on a measure the search optimises.
          \item \textbf{Architecture sweep.} A sweep over widths, depths and learning rates was programmed but no result was kept, so the network layers are those of the published design rather than a measured optimum.
          \item \textbf{Four implementation faults in the writing path}, none of which prevents a page being drawn correctly but each of which costs time or contradicts the documentation: the stroke ordering sorts rows by ascending $Y$ and so writes the bottom line first; the acceleration the generator models is not executed by the driver; the generator's work area of 200~mm exceeds the driver's measured travel of 194~mm in X; and because the driver has a single modal feed register, pen-up travel runs at the pen-down feed, lengthening every job.
          \item \textbf{Missing documentation and one inconsistent constant.} The CAD files, photographs, wiring diagram and camera data are outstanding, as is the number of candidate tries, which appears as 14, 20 and 30 in three places. \mustdo{reconcile the number of tries and supply the hardware documentation}
          \item \textbf{Training mode.} The proposal planned training on the single-board computer; training is done offline and the interface has no training function, because the board the budget allowed can run inference but cannot train these networks in a usable time.
          \item \textbf{The CTC loss is not safe against an empty transcription.} The gradient check found that a label sequence of length zero builds two shifted arrays at the wrong length and raises a shape error instead of returning a loss; every length of one or more passes, so the derivative is correct and this is a boundary indexing fault. It was left unfixed because no line in the data has an empty transcription, and a future data set would expose it.
        \end{enumerate}

        \textbf{What must change to meet R4, and what was measured rather than assumed.}
        Because the requirement fails on both halves, two independent changes are needed, and both were tested on the development computer so that the sizes quoted are measured factors and not estimates.
        The two candidate changes were applied to the same 14 pages, with the same weights, and the factor each produced was then applied to the stage times actually measured on the board.

        The first change is to execute the network forward pass differently.
        Loading the same weights into equivalent library modules and running every crop of a page in one batched forward pass, instead of the hand-written pass one crop at a time, is faster by a measured factor of 2.76 over 14 pages, with the text read out identical.
        Applied to the two on-target runs, that takes the network cost from 0.0602 to 0.0218~s per character and from 0.0394 to 0.0143~s per character, which is 0.9 and 0.6 times the $0.025\,$s per character term.
        \emph{The network half of R4 is therefore met by this change alone}, at both page lengths measured, and the change does not require a framework: batching the existing from-first-principles forward pass per page rather than per crop is the same arithmetic on a larger array.

        The second change must reduce the per-page image cost by about a factor of 8, from 16.1 to 16.8~s to under the fixed 2~s term, and the cheap way of doing it does not work.
        Running the image stages at half linear resolution reduced them by a measured factor of 4.37, which is most of what is needed, but the segmenter found 182 lines where it had found 247, not one of the twelve pages kept its line count, the text read back from the reduced page agreed with the full-resolution reading on 33.8\% of characters on average and 14.1\% at worst, and on two of the fourteen pages the segmenter failed outright.
        Decimation is therefore not an available remedy at any speed, and the reduction has to come from redesigning the line grouping, which is the part of the segmenter whose cost is not proportional to the pixel count.

        Taken together the two changes bring the projected totals to 14.6~s against an allowance of 14.50~s on the 500-character page and 17.9~s against 26.55~s on the 982-character page, that is, one just missed and one comfortably met.
        The honest conclusion is therefore not that R4 is unreachable but that it is reachable in two steps, of which the first is available today and measured, and the second requires work on the segmenter that this project did not do.
        It should be read alongside the economic constraint: a faster processing board would have reduced both terms at once without any of this work, and the budget did not allow one.

        \textbf{What must change for the other requirements.}
        For R2 on text, three things in order of expected effect: rebuild the training caches so that the corrected transcripts are actually used, which is a defect fix rather than an improvement; close the train-and-serve gap on the writer classifier and install the reader the measurements were taken with; and add a decoder that uses word context.
        The last would most likely close the 0.61-point gap on its own, given that character accuracy is 94.39\% while word accuracy is 73.47\%, but it must be implemented as a soft preference rather than a lexicon constraint, so that the arbitrary strings R1 requires remain writable.
        For R1, the character half needs the synthesiser to model letter context rather than compose isolated glyphs, which is the one change that would address the cursive writers where the whole 36-point spread lives; the style half needs a judge that took no part in the best-of-$N$ search, or a human panel, before the 85\% correlation can be claimed as a correlation.
        For R3, nothing is required to meet the specification, which is already met on the target platform; what is worth doing is measuring a block shorter than 43 characters to confirm where the margin actually closes, and correcting the modal feed register so that pen-up travel runs at the fast feed it was meant to.
        For R5, draw and measure at least five squares and five triangles with the diagonals, which would turn a compliance statement into a repeatability figure, and square the frame to remove the $0.706^{\circ}$ error the triangle identified.
        For R6, nothing; the requirement is met with one style in hand, and the margin would be restored by returning the two withdrawn writers once their pages are re-labelled.

    \subsubsection{Strong points of the design}
    \label{sec:dis-strong}

        Each strong point follows from a design decision.
        \begin{itemize}
          \item \textbf{The recogniser stands up on a public benchmark.} Reading whole lines removes the dependence on a character separation that cannot work for connected handwriting, and on 976 held-out crops of the public IAM corpus the from-first-principles reader reached 94.39\% character accuracy: 12.6 points above the published figure of Graves \emph{et al.} and 2.0 points below K{\i}z{\i}l{\i}rmak \emph{et al.} This is the one result in the report that can be compared with the literature on the same material, and it holds.
          \item \textbf{The conditioning generalises, and that is measured rather than asserted.} The flat-field correction, local thresholding and skew correction were designed so that an ordinary indoor photograph can be used. The evidence that they work is that the reader scores 94.39\% on crops cut by the corpus's distributor and 94.42\% on crops this project photographed with a webcam and cut with its own segmenter, a difference of 0.03 points. Two independent paths to the same figure is a stronger statement than either figure alone.
          \item \textbf{The writer classifier was chosen for the harder job, not the better score.} An ink-based classifier scored higher on photographs and was rejected because it cannot judge a constant-width plotter's output. The stroke-normalised classifier reaches 98.08\% per line and names all 11 writers correctly on held-out material, and it also works on the machine's own drawings, which is what makes the R1 measurement possible at all.
          \item \textbf{The gantry error is diagnosed, not just bounded.} The mean absolute error over a 100~mm square is 0.175~mm against a 0.5~mm limit, and the triangle identifies the dominant residual as a $0.706^{\circ}$ squareness error with residuals of 0.02~mm, at the caliper's resolution. The design therefore yields an actionable cause rather than a pass mark, and the absolute step bookkeeping prevents rounding from accumulating over a page.
          \item \textbf{Errors announce themselves.} Over the held-out lines the median confidence on a correct naming is 99.9\% and the mean on an incorrect one 55.8\%, so the classifier's own output separates its mistakes from its successes well enough to be used as a warning.
          \item \textbf{Provenance was treated as part of the measurement.} The transcript lookup that underpins qualification test~2 found that part of the material an earlier draft had reported as held out was in fact inside the reader's training corpus, and that several transcripts broke their lines in the wrong places. Correcting the data and discarding the contaminated material raised the reported character accuracy from 91.02\% to 94.39\% without touching a single weight. Finding that a figure was wrong in the favourable direction is the discipline working.
          \item \textbf{Safety in software:} end-stop confirmation over three consecutive samples, software limits, and automatic retraction after a trip.
          \item The algorithms are implemented from first principles and documented well enough to be analysed, and the known limitations are stated rather than left for a reader to find.
        \end{itemize}

    \subsubsection{Expected failure conditions}
    \label{sec:dis-failure}

        The system is expected to fail or degrade strongly under the following conditions.
        \begin{itemize}
          \item \textbf{Lighting.} Illuminance below 400~lux or strongly directional light with shadows violates FC1; the thresholding then merges ink with shadows or loses faint strokes, and recognition accuracy falls.
          \item \textbf{Page skew and perspective.} Large skew angles or a page that is not flat exceed the range of the skew correction.
          \item \textbf{Non-text regions.} Diagrams, tables and printed text can be segmented as text lines; the text versus non-text scoring reduces but does not exclude this, and the recogniser then produces meaningless characters.
          \item \textbf{Unseen writers.} The writer classifier is an eleven-class closed-set classifier, so a writer who is not among the eleven is assigned to one of them, because the softmax cannot express ``none of the above''.
          The confidence output mitigates this rather than solving it: over the held-out lines a correct naming carries a median confidence of 99.9\% while the three errors carried 39.5, 44.2 and 83.8\%, so a threshold would flag most such cases, but no threshold is applied and the system currently answers an unseen writer with a name and no warning.
          \item \textbf{Short reproductions.} R3 is met from 43 characters upward, and the margin closes as the block shortens because the best-of-$N$ search is paid once per block while the allowance is strictly proportional to length. From the measured synthesis cost and drawing rate the allowance is exhausted at about 35 characters, so a request of a word or two is expected to exceed the 4~s per character specification. This has not been measured, because no block that short was drawn.
          \item \textbf{The last line of a page is the weakest case for naming the writer.} All three naming errors on the held-out material fell on the last line of their page, which under a uniform hypothesis has probability $3.5\times10^{-4}$. A final line is usually a part line carrying a few words, so it offers the least ink to judge; the page-level decision was unaffected because pooling by ink weights such a line down, but a single short line submitted on its own is the input most likely to be misattributed.
          \item \textbf{Styles outside the glyph library.} A writer whose letter shapes differ greatly from the stored paths (unusual letters, capitals, digits or symbols that were scarce in the training pages) is reproduced with the nearest stored shape; characters never harvested cannot be written in that style.
          \item \textbf{Text with accents, mathematical notation or rare symbols}, which occur in the personal pages, are outside the benchmark vocabulary.
          \item \textbf{Motor stalls and lost steps.} The open-loop gantry has no feedback, so a stall shifts every later mark without any signal.
          \item \textbf{Limit switch failure.} A switch that does not close lets the carriage run against the mechanical stop; the software protection depends on the switches, and a hardware emergency stop that cuts the power does not exist.
          \item \textbf{Long sentences.} The synthesis time grows with the number of characters and tries, so a long text at a high quality setting would take minutes on the board and may time out the web request.
          \item \textbf{Network failure} between the browser and the server prevents control of the system; a running job continues on the board.
        \end{itemize}
        The system does not fail gracefully in all cases, and the two kinds of failure are not equally fixable.
        The classifier returns an answer for an unseen writer without any warning, but it does carry the information needed to warn: its confidence separates its own errors from its successes on the held-out material, so a threshold would convert a silent wrong answer into a refusal.
        That is an afternoon's work and was not done.
        The gantry, by contrast, cannot detect lost steps at all, because there is no encoder and the budget did not allow one; a stall shifts every later mark on the page with no signal of any kind, and no amount of software can recover it.
        The first failure mode is therefore a missing feature and the second is a consequence of the design.

\subsection{Considerations in the design} % REQUIRED HEADING
\label{sec:dis-considerations}

This subsection addresses the ergonomics, safety, environmental, social, legal and ethical aspects of the design.

    \subsubsection{Ergonomics}

        The user works through a web interface of four pages (Read, Write, Model and Style) in a browser on any laptop on the network, so nothing is installed.
        That the interface is browser-based is not only a convenience: the processing board has no HDMI output, so it has no local display at all, and a networked interface was the only way to operate the machine.
        What began as a constraint turned out to be the better ergonomic choice, because the operator works from a full keyboard and screen beside the gantry rather than from a console attached to it.
        The Read page captures the page with a live preview and shows the recognised text, the writer, the confidence and the measured time against the specification's allowance; the Write page keeps the writing grid locked until the gantry is calibrated and shows a G-code preview before sending.
        The shortcomings are that the live progress of a running job is not displayed and the interface has no training function, the latter because the board cannot train these networks in a usable time.
        The rig is set up by hand, with the page, the pen and the paper placed by hand; its documentation is incomplete.

    \subsubsection{Health and safety}

        The electronics operate at low voltage: a single-board computer, stepper-driver modules and small motors, with the driver supply run at 24~V, so the risk of electric shock is low. \mustdo{the mains insulation rating of the power supply, which is a property of the bought unit and was not recorded.}
        The mechanical hazards are the moving carriage, belts and pulleys, which are pinch points while the gantry moves; the pen tip is not hazardous at the writing force.
        The protective measures are the end-stop switches with confirmation over three consecutive samples, the software limits, the automatic retraction and calibration at a low speed (7.8~mm/s).
        A running job can also be stopped from the interface: the server exposes a stop endpoint, the motion loop tests an abort flag between steps, and the recorded test log shows a job being stopped, the pen lifted and the carriage parked at the corner in response.
        What does not exist is a hardware emergency stop that removes motor power independently of the software, so every protective measure depends on the program continuing to run; fitting a mains-independent emergency stop is recommended before the machine is used by anyone other than its designer.
        Users must keep their hands away from the carriage while a job runs.

    \subsubsection{Environmental impact}

        The ODROID N2+ is a low-power board and the stepper motors draw power only while the gantry works, so the energy use is small beside a desktop computer.
        The board was also obtained second hand rather than bought new, which was a budget decision before it was an environmental one but has the same effect: the most complex electronic component in the system is a reused one, and the embodied manufacturing cost of a new single-board computer was avoided entirely.
        The system uses paper; test pages and failed jobs add waste, limited by the preview before writing and by the option to write on used paper.
        The mechanical parts printed in plastic or bought as standard parts are \mustdo{which parts were 3D-printed and with which material}; they can be reused, and plastic and electronic waste must go to a recycling point, not into domestic waste.
        The pen ink is a small consumable, no hazardous material beyond the electronics is used, and the only emissions are the motor noise and the pulses of the drive.

    \subsubsection{Social and legal impact}

        The social benefit is that the system can write in a person's own handwriting for those who cannot write neatly or at all, for example after an injury, and offers the personal touch the proposal identified.
        There is also a risk of misuse: reproducing a person's handwriting can serve forgery or impersonation, for instance on signed documents or letters.
        The plotter pen's constant line width reduces but does not remove the risk, because it differs from natural writing in stroke width and pressure (the ink-based classifier identified only about 20\% of the synthesised output, so these ink properties are not reproduced), yet a document in a plausible style can still mislead a reader.
        Handwriting samples are personal information and a writer profile is similar to biometric data, since it permits identification of the writer; in South Africa the Protection of Personal Information Act (POPIA) applies to such data, requiring the consent of the writers, the minimum storage and the protection of the profiles against access by others.
        The samples of the five personal writers require documented consent, \mustdo{confirm that consent was given and in what form}, and the eight database writers contributed to a public research database whose terms of use apply.
        The server has no user accounts (only an optional shared key), accepts requests from any origin and can start physical motion for anyone who can reach its port, so it must be operated only on a trusted local network.
        Beyond the general electrical safety rules no other legislation was identified for a research prototype of this type; \mustdo{any applicable regulation or university rule}.

    \subsubsection{Ethics clearance}

        The project used handwriting samples of five named persons and of a public database, so it involves human participants in a limited way.
        The status of ethics clearance is \mustdo{ethics clearance number, or a statement that no clearance was required}.
        If the human panel of qualification test~1 is carried out, the raters give informed consent and their responses are stored without personal identifiers, and clearance must be confirmed before the panel is run.

